% Generated draft: requires mathematical and source review.
\documentclass[11pt]{article}
\usepackage[margin=25mm]{geometry}
\usepackage{fontspec}
\setmainfont{Latin Modern Roman}
\usepackage{amsmath,amssymb,mathtools,unicode-math}
\setmathfont{Latin Modern Math}
\usepackage{enumitem,needspace,xurl,hyperref}
\setlength{\parindent}{0pt}
\setlength{\parskip}{5pt}
\setlength{\emergencystretch}{3em}
\newcommand{\N}{\mathbb{N}}
\newcommand{\Z}{\mathbb{Z}}
\newcommand{\Q}{\mathbb{Q}}
\newcommand{\R}{\mathbb{R}}
\newcommand{\C}{\mathbb{C}}
\newcommand{\F}{\mathbb{F}}
\newcommand{\A}{\mathbb{A}}
\newcommand{\PP}{\mathbb{P}}
\newcommand{\E}{\mathbb{E}}
\newcommand{\eps}{\varepsilon}
\newcommand{\dd}{\,\mathrm{d}}
\begin{document}
\section*{UnsolvedMath ID: 1900001}
\textbf{Source problem code:} catalogue entry
\subsection{Definitions and mathematical statement}
Karpenkov's question is to find an integer cosine rule for integer triangles in integer trigonometry. In the usual lattice-geometric reading, an integer triangle is a nondegenerate triangle with vertices $A,B,C\in\Z^2$, so
\[
\det(B-A,C-A)\ne0.
\]
This is a lattice-vertex condition, not the Euclidean condition that all three ordinary side lengths are integers. However, the supplied source extract itself does not define ``integer triangle'', so this reading remains a convention to confirm against the cited source.

The requested rule is to be an analogue of the cosine rule expressed in the intended integer-trigonometric side and angle invariants and valid for every admissible integer triangle, with all necessary qualifications. No definition of integer length, integer cosine, integer sine, angle equivalence, orientation, or the allowed data for such a rule is included. Consequently the record does not determine an explicit identity or a unique quantified conjecture that can be stated without importing absent definitions. In particular, replacing the requested invariants by ordinary Euclidean lengths and cosine would change the subject rather than formalize the intended integer rule.

The supplied research summary labels the question open with moderate confidence and says it located no verified citation for a complete rule for arbitrary integer triangles. This is a supplied source-search limitation, not independent evidence that no rule exists. The title also mentions the ``IKEA problem'', but no separate IKEA statement is supplied and none is reconstructed here.
\subsection{Short English statement}
Find a cosine-rule analogue in Karpenkov's integer-trigonometry framework for all integer triangles. The extract omits the definitions of the relevant lattice lengths and angle invariants, so it cannot yet specify the desired identity. Its supplied open-status assessment is tentative and reports no verified settling reference.
\subsection{Sources}
\begin{itemize}
\item Local source registry: \url{data/unsolvedmath/problems.json}, permanent UnsolvedMath ID $1900001$, source code catalogue entry, including the supplied OPEN-TRIAGE summary. Dataset attribution: Ulam AI and UnsolvedMath Contributors; original-source attribution and rights are retained.
\item Karpenkov, \emph{Open problems in geometry of continued fractions} (2017), Problem 1, as identified in the source list: \url{https://arxiv.org/abs/1712.01450}. The metadata reports source-TeX extraction and access on July 29, 2026, and flags current status and rights for review.
\item The supplied research summary reports no verified citation settling the complete integer cosine-rule question. No independent source check or inspection of the linked document was performed for this draft.
\end{itemize}
\end{document}
